\begin{frame}[label=q5c3]{5c.3 · DataFrame: RIGHT OUTER JOIN}
  \small
  \pregunta{¿Hay meses sin registros entre enero de 2019 y diciembre de 2024?}
  \entrada{Agregados por \texttt{anio, mes} a la izquierda; calendario independiente de 72 meses a la derecha.}
  \flujo{\texttt{groupBy(anio,mes)}}{\texttt{RIGHT OUTER JOIN}\\Conservar calendario}{\texttt{coalesce} a cero\\Filtrar total = 0}{72 meses presentes\\0 meses sin registros}
  \begin{columns}[T,onlytextwidth]
    \column{0.55\textwidth}
    \textbf{Código Scala original (fragmento)}
    \codigo{sql/05c_3_right_calendario}{16}{18}
    \codigo{sql/05c_3_right_calendario}{21}{21}
    \column{0.42\textwidth}
    \textbf{Salida}\\[5pt]
    {\footnotesize
\begin{tabular}{@{}ll@{}}
\toprule
Control & Resultado\\
\midrule
Meses conservados & 72\\
Meses sin registros & 0\\
Personas & 121\,427\\
\bottomrule
\end{tabular}
    }
  \end{columns}
  \respuesta{Todos los meses del calendario tienen registros. A diferencia del catálogo geográfico, el calendario se genera sin depender de las filas observadas.}
  \origen{sql}{05c_3_right_calendario}
\end{frame}
